\phantomsection
\addcontentsline{toc}{section}{Introduction}

% Définition de la commande ui pour les interfaces
\newcommand{\ui}[4]{%
  \begin{minipage}[t]{0.48\linewidth}
    \centering
    \includegraphics[width=\linewidth]{images/#2}%
    \caption{#1}%
    \parbox{\linewidth}{\small #3}%
    \parbox{\linewidth}{\footnotesize\textit{#4}}%
  \end{minipage}%
}

\section*{Introduction}
Dans ce chapitre, nous présentons d'abord les interfaces majeures de l'application
Inter\,Actif et de la documentation de l'API. Ces interfaces constituent
les principaux écrans et vues avec lesquels un organisateur ou un participant
interagit. Nous décrivons ensuite les limitations observées lors de la mise
en œuvre, afin de fournir un aperçu complet de la solution proposée.

\section{Interfaces de documentation de l'API et de l'application}
\subsection{Interfaces de l'application}

% LIGNE 1 : Accueil et Connexion
\begin{figure}[htbp]
  \centering
  \ui{Accueil}{page_acceul.jpg}%
     {Écran d'entrée présentant les fonctionnalités clés (sondages, quiz, chat) et deux parcours : Organisateur ou Rejoindre.}%
     {Orienter immédiatement l'utilisateur vers l'action adaptée.}
  \hfill
  \ui{Connexion sécurisée}{page_connexion.jpg}%
     {Formulaire e‑mail et mot de passe avec liens « S'inscrire » et « Mot de passe oublié ».}%
     {Authentifier l'utilisateur pour un accès personnalisé.}
\end{figure}

% LIGNE 2 : Inscription et Réinitialisation
\begin{figure}[htbp]
  \centering
  \ui{Inscription sécurisée}{page_inscription.jpg}%
     {Création de compte : nom complet, e‑mail, mot de passe (avec indicateur de robustesse) et confirmation.}%
     {Enrôler de nouveaux utilisateurs en garantissant des mots de passe forts.}
  \hfill
  \ui{Réinitialisation du mot de passe}{page_mise_ajour.jpg}%
     {Saisie de l'adresse e‑mail pour recevoir un lien de réinitialisation.}%
     {Rétablir l'accès en cas d'oubli de mot de passe.}
\end{figure}

% LIGNE 3 : Tableau de bord et Mes événements
\begin{figure}[htbp]
  \centering
  \ui{Tableau de bord}{page_dashboard.jpg}%
     {Carte de bienvenue personnalisée et statistiques (événements, participants, en cours/à venir) avec action « Nouvel événement ».}%
     {Donner une vue d'ensemble et un accès rapide aux actions clés.}
  \hfill
  \ui{Mes événements}{page_evernement.jpg}%
     {Liste des événements avec recherche, filtres (\emph{Tous}, \emph{Actifs}, \emph{À venir}) et bouton d'ajout.}%
     {Retrouver, filtrer et créer des événements aisément.}
\end{figure}

% LIGNE 4 : Détail d'un événement et Création d'un événement
\begin{figure}[htbp]
  \centering
  \ui{Détail d'un événement}{page_dtail_event.jpg}%
     {Récapitulatif (nom, date/heure, code d'accès, description) et sections \emph{Sondage}, \emph{Quiz}, \emph{Chat} avec compteurs et actions.}%
     {Centraliser la gestion des modules d'interaction.}
  \hfill
  \ui{Création d'un événement}{page_creation_event.jpg}%
     {Formulaire par étapes : informations générales, date/heure (début et fin optionnelle), lieu et participation.}%
     {Saisir toutes les métadonnées nécessaires à la planification.}
\end{figure}

% LIGNE 5 : Liste des modules et Création de module Quiz
\begin{figure}[htbp]
  \centering
  \ui{Modules — Quiz (liste)}{page_des_moduls_quiz.jpg}%
     {Recherche, filtre, bouton « + » et cartes modules (état actif, nombre de questions, action « Voir »).}%
     {Gérer les différents modules de quiz d'un événement.}
  \hfill
  \ui{Création d'un module — Quiz}{page_creation_module.jpg}%
     {Titre du module, type de réponse (choix unique/multiple, échelle, texte libre), éditeur de questions et options.}%
     {Concevoir des activités interactives adaptées aux objectifs pédagogiques.}
\end{figure}

% LIGNE 6 : Création de module Sondage et Participation (choix unique)
\begin{figure}[htbp]
  \centering
  \ui{Création d'un module — Sondage}{Screenshot_20250906-164855.jpg}%
     {Interface de création de sondage avec les mêmes types de réponses que le quiz, orientée recueil d'opinions.}%
     {Construire rapidement des sondages ciblés.}
  \hfill
  \ui{Participation — Choix unique}{page_participation.jpg}%
     {Affichage d'une question à choix unique avec boutons radio et navigation \emph{Précédent/Suivant}.}%
     {Permettre une réponse rapide et fluide des participants.}
\end{figure}

% LIGNE 7 : Participation (choix multiples) et Résultats de question
\begin{figure}[htbp]
  \centering
  \ui{Participation — Choix multiples}{Screenshot_20250906-165209.jpg}%
     {Question à choix multiples avec cases à cocher permettant plusieurs réponses.}%
     {Recueillir des réponses multiples lorsque plusieurs solutions sont pertinentes.}
  \hfill
  \ui{Résultats d'une question de quiz}{page_module_detail.jpg}%
     {Synthèse : nombre de répondants, barres de pourcentage par option, navigation entre les questions.}%
     {Visualiser instantanément la compréhension et l'adhésion du public.}
\end{figure}

% LIGNE 8 : Participants et Profil utilisateur
\begin{figure}[htbp]
  \centering
  \ui{Participants}{page_participants.jpg}%
     {Statistiques (participants, réponses, moyenne, événements, actifs/à venir), recherche et liste détaillée ; action \emph{Inviter}.}%
     {Suivre l'engagement et administrer la base de participants.}
  \hfill
  \ui{Profil utilisateur}{page_profil.jpg}%
     {Nom, e‑mail et rôle, possibilité d'édition via un crayon et bouton \emph{Déconnexion}.}%
     {Consulter et gérer son identité et ses droits.}
\end{figure}

% LIGNE 9 : Paramètres du compte (réinitialisation)
\begin{figure}[htbp]
  \centering
  \ui{Paramètres du compte}{page_paramtre.jpg}%
     {Sélection des catégories de données à effacer (événements, participants, modèles, données locales) puis exécution de la réinitialisation.}%
     {Offrir un contrôle explicite sur la purge des données et la remise à zéro.}
  \hfill
  \begin{minipage}[t]{0.48\linewidth}
    % Espace réservé pour une interface supplémentaire éventuelle.
  \end{minipage}
\end{figure}

\subsection{Interfaces de documentation de l'API}
\textbf{Remarque.} Les captures d'écran de la documentation de l'API ne sont pas
fournies. La documentation se présente généralement sous forme de pages Web
interactives (Swagger ou Redoc) où l'on retrouve :
\begin{itemize}
  \item \textbf{Aperçu et authentification :} description de l'URL de base,
    des schémas d'authentification (clé API ou Bearer) et des en-têtes
    nécessaires pour chaque requête.
  \item \textbf{Ressources principales :} endpoints REST tels que
    \texttt{/events}, \texttt{/modules}, \texttt{/participants}, \texttt{/responses},
    chacun documenté avec ses méthodes autorisées (GET, POST, PATCH, DELETE)
    et les codes de réponse associés.
  \item \textbf{Schémas JSON :} modèles détaillés pour les requêtes et
    réponses, précisant les champs obligatoires et optionnels, avec des
    exemples de valeurs.
  \item \textbf{Exemples d'utilisation :} requêtes cURL ou extraits de
    code illustrant la création d'un événement, l'ajout d'un module, la
    soumission d'une réponse, ou la mise à jour de ressources existantes.
  \item \textbf{Gestion des erreurs :} format de retour des erreurs (code
    et message) et description des principaux cas d'échec (identifiants
    invalides, données manquantes, etc.).
\end{itemize}

\section{Limites de l'application}
\begin{itemize}[leftmargin=1.2em]
  \item \textbf{Dépendance réseau :} l'expérience en temps réel (quiz,
    sondages et chat) requiert une connexion stable ; de mauvaises
    conditions réseau peuvent altérer la réactivité et l'affichage des
    résultats.
  \item \textbf{Montée en charge :} un grand nombre de participants
    simultanés peut nécessiter un dimensionnement serveur ou l'ajout de
    mécanismes de pagination et de files d'attente pour éviter une surcharge.
  \item \textbf{Hétérogénéité des appareils :} la diversité des tailles
    d'écran et des versions d'OS implique d'adapter l'interface (polices,
    densité, accessibilité) afin d'assurer une expérience cohérente.
  \item \textbf{Confidentialité :} la gestion des données personnelles
    impose de respecter des exigences strictes (conformité RGPD) pour
    l'authentification, la conservation et la suppression définitive.
  \item \textbf{Fonctionnalités hors ligne limitées :} certaines actions
    critiques (vote en temps réel, synchronisation des résultats) ne sont
    pas disponibles sans accès internet.
\end{itemize}

\phantomsection
\addcontentsline{toc}{section}{Conclusion}
\section*{Conclusion}
Dans ce chapitre, nous avons décrit en détail les interfaces de
l'application Inter\,Actif en les illustrant de captures d'écran et en
explicitant leur objectif. Nous avons également expliqué la structure
générale de la documentation de l'API et souligné les limites inhérentes
à l'implémentation actuelle. Ces éléments constituent un socle utile pour
comprendre le fonctionnement global de la solution et orienter les
améliorations futures.
